% ============================================================
% 前置部分：摘要、关键词、致谢
% ============================================================

% ------------------------------------------------------------
% 中文摘要
% ------------------------------------------------------------
\chapter*{摘\ \ 要}
\addcontentsline{toc}{chapter}{摘要}

存内计算（Compute-in-Memory, CIM）架构通过在存储阵列内部直接执行矩阵--向量乘法，消除了制约传统数字系统的数据搬运瓶颈。
在诸多新兴器件路线中，有机光电子器件凭借其多级电导可调性、低压操作与本征光敏性，为神经形态视觉系统提供了独特的物理载体。
然而，将器件级进展转化为系统级视觉性能，仍面临器件变异性、有限模数转换器（Analog-to-Digital Converter, ADC）精度、保留漂移及跨制造实例可迁移性等关键挑战。

本研究建立了一个面向有机光电子存内计算阵列的可复用硬件感知训练（Hardware-Aware Training, HAT）仿真框架，
支持具有真实器件非理想性的视觉 Transformer（Vision Transformer, ViT）与卷积神经网络（Convolutional Neural Network, CNN）。
该框架采用配置文件驱动设计，可将文献报道的器件特性（电导窗口、周期到周期变异性、器件到器件失配、保留动力学、光响应非线性、写入非线性）映射为任务级视觉精度。

实验首先系统评估了 ResNet-18、ConvNeXt-Tiny 与 Tiny-ViT-5M 在 CIFAR-10、CIFAR-100、Flowers-102 及 ImageNet-1k 上的 baseline 精度。
在此基础上，本研究提出 HAT 训练范式的系统分类学，并发现\emph{硬件实例过拟合}现象：
标准固定掩膜 HAT 在训练时硬件实例上表现良好，但在全新器件失配实例上崩溃至接近随机水平（CIFAR-10 上 $10.00\pm0.00\%$）。
为应对此问题，本研究提出 Ensemble HAT——在每个训练轮次边界重采样结构化器件到器件（Device-to-Device, D2D）失配场——
将 fresh-instance 精度恢复至 $86.16\pm0.19\%$。

进一步地，本研究构建了面向有机光电子 CIM 部署的\emph{失效模式图谱}（failure mode atlas），
系统刻画了多种可复现的失效模式及其边界：
（1）硬件实例过拟合与 Ensemble HAT 缓解；
（2）严重写入非线性（$NL=2.0$）下的双注意力通路崩溃；
（3）多层感知机（Multi-Layer Perceptron, MLP）线性补偿的 $32\%$ fresh-instance 诊断上界；
（4）空间相关 D2D 的有界退化；
（5）保留漂移的 $79\%$ 渐近平台期。
所有失效模式均通过严格的负结果（negative results）加以界定，从而将瓶颈定位为结构性泛化障碍而非优化缺口。

本研究的贡献包括：
（1）一个可复用、可扩展的有机光电子 CIM 仿真框架；
（2）HAT 范式的系统分类学与硬件实例过拟合的发现；
（3）一个受经验数据约束的部署包络，明确了哪些噪声模型与训练干预组合能够迁移至新硬件实例、哪些不能。
这些结果为有机光电子神经形态硬件的算法--器件协同设计提供了系统级决策依据。

\vspace{1em}
\noindent\textbf{关键词：}存内计算；硬件感知训练；视觉 Transformer；有机光电子器件；硬件实例过拟合

\vspace{2em}

% ------------------------------------------------------------
% 英文摘要
% ------------------------------------------------------------
\chapter*{Abstract}
\addcontentsline{toc}{chapter}{Abstract}

Compute-in-memory (CIM) architectures eliminate the data-movement bottleneck of conventional digital systems by performing matrix--vector multiplications directly inside memory arrays.
Among emerging device routes, organic optoelectronic devices offer a distinctive physical substrate for neuromorphic vision systems thanks to their multi-level conductance tunability, low-voltage operation, and intrinsic photosensitivity.
Yet translating device-level progress into system-level visual performance remains hindered by device variability, limited analog-to-digital converter (ADC) resolution, retention drift, and cross-fabrication-instance transferability.

This thesis presents a reusable hardware-aware training (HAT) simulation framework for organic optoelectronic CIM arrays that supports vision transformers (ViT) and convolutional neural networks (CNN) under realistic device non-idealities.
The framework adopts a profile-driven design that maps literature-reported device characteristics---conductance window, cycle-to-cycle variability, device-to-device (D2D) mismatch, retention dynamics, photoresponse nonlinearity, and write nonlinearity---into task-level visual accuracy.

Experiments first systematically evaluate ResNet-18, ConvNeXt-Tiny, and Tiny-ViT-5M on CIFAR-10, CIFAR-100, Flowers-102, and ImageNet-1k.
Building on these baselines, the thesis introduces a taxonomy of HAT regimes and discovers \emph{hardware-instance overfitting}:
standard fixed-mask HAT performs well on the training-time hardware instance yet collapses to near-chance level on fresh mismatch realizations ($10.00\pm0.00\%$ on CIFAR-10).
To address this, the thesis proposes Ensemble HAT, which resamples structured D2D mismatch fields at epoch boundaries during training, recovering fresh-instance accuracy to $86.16\pm0.19\%$.

Furthermore, the thesis constructs a \emph{failure mode atlas} for organic optoelectronic CIM deployment,
systematically cataloguing multiple reproducible breakdown patterns and their boundaries:
(1) hardware-instance overfitting and its Ensemble HAT remedy;
(2) dual-attention-pathway collapse under severe write nonlinearity ($NL=2.0$);
(3) the $32\%$ fresh-instance diagnostic ceiling of MLP-only linearization;
(4) bounded degradation under spatially correlated D2D;
(5) the $79\%$ asymptotic retention plateau.
Each failure mode is bounded by rigorous negative results, locating bottlenecks as structural generalization barriers rather than optimization gaps.

Contributions include:
(1) a reusable, extensible simulation framework for organic optoelectronic CIM;
(2) a taxonomy of HAT regimes and the discovery of hardware-instance overfitting;
(3) an empirically constrained deployment envelope that delineates which noise-model and training-intervention combinations transfer to new hardware instances and which do not.
These results provide system-level decision support for algorithm--device co-design of organic optoelectronic neuromorphic hardware.

\vspace{1em}
\noindent\textbf{Keywords:} compute-in-memory; hardware-aware training; vision transformer; organic optoelectronic device; hardware-instance overfitting

\vspace{2em}

% ------------------------------------------------------------
% 致谢
% ------------------------------------------------------------
\chapter*{致\ \ 谢}
\addcontentsline{toc}{chapter}{致谢}

本研究在课题组的持续讨论、计算资源支持与多轮论文审阅中逐步完成。

衷心感谢导师在整个研究过程中给予的悉心指导与学术启发。
感谢实验室同门在器件表征、仿真调试与论文写作中提供的宝贵建议与无私帮助。
感谢相关合作团队在有机光电子器件制备、测量与建模讨论方面提供的支持。

最后，感谢家人在求学期间给予的持续理解与鼓励。

\vspace{2em}
\noindent
{李松桥}\\
{2026年5月}
